\subsection{The intersection of a finite number of valuation rings}

PROPOSITION 1. Let K be a field, \((\mathbf{A}_{i})_{1 \leqslant i \leqslant n}\) a finite family of valuation rings of K and \(B = \bigcap_{i=1}^{n} A_{i}\) . We write \(p_{i} = B \cap m(A_{i})\) . Then \(A_{i} = B_{p_{i}}\) for all i and the field of fractions of B is K.

Clearly \(B_{p_t} \subset A_t\) . To prove the converse inclusion we need the following lemma: LEMMA 1. Let \(v_i (1 \leqslant i \leqslant n)\) be valuations on the field \(K\) and \(x \in K^*\) . Then there exists a polynomial \(f(X)\) of the form

\[
\begin{array}{l} (1) f (\mathbf {X}) = 1 + n _ {1} \mathbf {X} + \dots + n _ {k - 1} \mathbf {X} ^ {k - 1} + \mathbf {X} ^ {k} \\ \quad (k \geqslant 2, n _ {j} \in \mathbf {Z} for 1 \leqslant j \leqslant k - 1) \end{array}
\]

such that \(f(x) \neq 0\) and the element \(z = f(x)^{-1}\) enjoys the following properties for \(1 \leqslant i \leqslant n\) :

\[
\begin{array}{l l} v _ {i} (z) = 0 & \text { if } \quad v _ {i} (x) \geqslant 0 \\ v _ {i} (z) + v _ {i} (x) > 0 & \text { if } \quad v _ {i} (x) <   0. \end{array}
\]

Assuming this lemma for a moment, we show how it implies that \(A_{1} \subset B_{p_{1}}\) . Let x be a non-zero element of \(A_{1}\) . We apply the lemma to x and valuations \(v_{i}\) associated with the \(A_{i}\) . Then \(v_{i}(z) \geqslant 0\) and \(v_{i}(zx) \geqslant 0\) for all i, hence \(z \in B\) and \(zx \in B\) . As \(v_{1}(x) \geqslant 0\) , \(v_{1}(z) = 0\) and hence \(z \notin p_{1}\) . Hence \(x = xz/z \in B_{p_{1}}\) . The field of fractions of B then contains \(A_{1}\) and hence is K.

We now pass to the proof of the lemma. Let I be the set of indices \(i\) such that \(v_{i}(x) \geqslant 0\) . For all \(i \in I\) , let \(\bar{x}_{i}\) denote the canonical image of \(x\) in \(\kappa(A_{i})\) . For all \(i \in I\) we construct a polynomial \(f_{i}\) as follows: if there exists a polynomial \(g(X)\) of the form (1) such that \(g(\bar{x}_{i}) = 0\) in \(\kappa(A_{i})\) , we take \(f_{i}\) to be such a polynomial; otherwise we take \(f_{i} = 1\) . Then we write \(f(X) = 1 + X^{2} \prod_{i \in I} f_{i}(X)\) . It is obviously a polynomial of the form (1). If \(i \in I\) , then \(f(x) \in A_{i}\) and also \(f(\bar{x}_{i}) \neq 0\) by construction; hence \(f(x) \notin m(A_{i})\) , \(v_{i}(f(x)) = 0\) and \(v_{i}(z) = 0\) . If \(i \notin I\) , then \(v_{i}(x) < 0\) , whence \(v_{i}(f(x)) = kv_{i}(x)\) ( \(\S\) 3, no. 1, Proposition 1) and

\[
v _ {i} (x) + v _ {i} (z) = (1 - k) v _ {i} (x) > 0
\]

(since \(k \geqslant 2\) ). Whence the lemma.

PROPOSITION 2. With the hypotheses of Proposition 1 suppose further that \(\mathbf{A}_i \not\subset \mathbf{A}_j\) for \(i \neq j\) . Then the \(\mathfrak{p}_i\) are distinct maximal ideals of \(\mathbf{B}\) and every maximal ideal of \(\mathbf{B}\) is equal to one of the \(\mathfrak{p}_i\) .

If \(\mathfrak{p}_i\subset \mathfrak{p}_j\) for \(i\neq j,\mathrm{A}_i = \mathrm{B}_{\mathfrak{p}_i}\supset \mathrm{B}_{\mathfrak{p}_j} = \mathrm{A}_j\) . It is then sufficient to apply Chapter II, § 3, no. 5, Corollary to Proposition 17.

COROLLARY 1. Suppose that \(\mathbf{A}_i \subsetneq \mathbf{A}_j\) for \(i \neq j\) . For every family of elements \(a_i \in \mathbf{A}_i\) ( \(1 \leqslant i \leqslant n\) ), there exists \(x \in \mathbf{B}\) such that \(x \equiv a_i \pmod{\mathfrak{m}(\mathbf{A}_i)}\) for \(1 \leqslant i \leqslant n\) .

Since the \(p_i\) are maximal ideals of B, \(A_i / m(A_i) = B_{p_i} / p_i B_{p_i} = B / p_i\) and it may therefore be assumed that \(a_i \in B\) for all \(i\) . The corollary then follows from the

fact that the canonical mapping from B to \(\prod_{i=1}^{n} (B/p_i)\) is surjective (Chapter II, § 1, no. 2, Proposition 5).

COROLLARY 2. Suppose that \(\mathbf{A}_i \not\subset \mathbf{A}_j\) for \(i \neq j\) . There exist elements \(x_i\) ( \(1 \leqslant i \leqslant n\) ) of \(\mathbf{K}\) such that \(v_i(x_i) = 0\) and \(v_j(x_i) > 0\) for \(i \neq j\) .

For each index \(i\) apply Corollary 1 to the family \((a_{j})\) such that \(a_{i} = 1\) and \(a_{j} = 0\) for \(j \neq i\) .

COROLLARY 3. Every valuation ring of K containing B contains one of the A \(_{i}\) .

We may confine our attention to the case where \(A_{i} \not\subset A_{j}\) for \(i \neq j\) . Let V be a valuation ring of K containing B. We write

\[
\mathfrak {p} = \mathfrak {m} (\mathrm{V}) \cap \mathrm{B}.
\]

There exists a maximal ideal \(p_{i}\) of B containing p, whence

\[
\mathrm{A} _ {i} = \mathrm{B} _ {\mathfrak {p} _ {i}} \subset \mathrm{B} _ {\mathfrak {p}} \subset \mathrm{V}.
\]

